% =============================================================================
%  Notas del Profesor — Sesión 02: Introducción a C++
%  Programación Avanzada — Universidad Panamericana
%  Compilar: pdflatex sesion02_introduccion_cpp_profesor.tex
% =============================================================================
\documentclass[12pt, oneside]{article}
\usepackage[profesor]{progavanzada}   % versión profesor (portada azul UP)

\begin{document}

\portada{Sesión 02}{Introducción a C++}{2026}

\tableofcontents
\newpage

% =============================================================================
\section{Objetivos de la sesión}
% =============================================================================

Al finalizar la sesión el alumno será capaz de:

\begin{enumerate}
  \item Explicar qué es C++ y en qué contextos se usa.
  \item Identificar y acceder a las fuentes de documentación principales.
  \item Distinguir entre compiladores e intérpretes, y nombrar los más usados.
  \item Compilar y ejecutar un programa simple desde la consola con \cpp{g++}.
  \item Comprender la estructura mínima de un programa C++ y el rol de cada línea.
\end{enumerate}

\begin{notaprofesor}[Duración estimada]
  50 minutos. Sugerido: 15 min expositivo + 20 min demostración en vivo
  (compilar en consola y abrir el cuaderno) + 15 min práctica guiada.
\end{notaprofesor}

% =============================================================================
\section{Secuencia de clase}
% =============================================================================

\subsection{Apertura — ¿Por qué C++? (10 min)}

Abrir con la pregunta: \textit{``¿Qué lenguajes ya conocen?''}.
Conectar C++ con aplicaciones que los alumnos usan (Chrome, videojuegos,
sistemas embebidos). Mostrar la tabla de características sin entrar en detalle.

\begin{notaprofesor}
  Si el grupo es de primer semestre de ingeniería, enfatizar el acceso a memoria
  como diferenciador clave frente a Python o Java. Si ya vienen con experiencia,
  enfocarse en el modelo de compilación y tipado estático.
\end{notaprofesor}

\subsection{Documentación (5 min)}

Abrir \href{https://en.cppreference.com}{cppreference.com} en el navegador y
buscar en vivo \cpp{std::cout}. Mostrar la estructura de la página: descripción,
parámetros, ejemplo, notas. El objetivo es que pierdan el miedo a la documentación
técnica desde el primer día.

\subsection{Compiladores y entornos (10 min)}

Explicar la diferencia \textbf{compilar vs interpretar} con el diagrama:

\begin{center}
  \texttt{main.cpp} $\xrightarrow{\text{g++}}$ \texttt{hola} (binario) $\xrightarrow{}$ ejecuta
\end{center}

Contrasta con cling/xeus-cling: evalúa línea a línea, sin generar binario.

\begin{notaprofesor}[Pregunta de sondeo]
  ``Si cling es más cómodo, ¿por qué no se usa en producción?''\\
  Respuesta esperada: rendimiento, dependencias, no genera binario distribuible.
\end{notaprofesor}

\subsection{Demostración en consola (15 min)}

Abrir terminal. Ejecutar en vivo:

\begin{codigo}[language=bash, numbers=none]
g++ --version
g++ -std=c++17 main.cpp -o hola
./hola
\end{codigo}

Luego introducir un error a propósito (quitar el punto y coma) y mostrar cómo
leer el mensaje de error: \texttt{archivo:línea:columna: error: descripción}.

\subsection{Cuaderno Jupyter (10 min)}

Abrir \texttt{sesion02\_introduccion\_cpp.ipynb}. Ejecutar las celdas
secuencialmente. Señalar:
\begin{itemize}
  \item No se necesita \cpp{main()}.
  \item Los \cpp{\#include} persisten entre celdas.
  \item \cpp{cin} no funciona interactivamente — mostrar la celda adaptada.
\end{itemize}

% =============================================================================
\section{Preguntas frecuentes}
% =============================================================================

\begin{notaprofesor}[FAQ]
\textbf{``¿Por qué \cpp{using namespace std;}? ¿Es mala práctica?''}\\
En proyectos grandes sí: puede causar colisiones de nombres. En código de clase
y archivos pequeños es aceptable para reducir ruido visual. Mencionarlo desde
el principio forma buenos hábitos.

\vspace{0.5em}
\textbf{``¿Qué pasa si omito \cpp{return 0;} en \cpp{main}?''}\\
Desde C++11 es opcional en \cpp{main}. El estándar garantiza retorno 0 implícito.
En otras funciones con tipo de retorno distinto de \cpp{void}, sí es obligatorio.

\vspace{0.5em}
\textbf{``¿\cpp{endl} o \cpp{\\n}?''}\\
\cpp{\\n} es más rápido (no hace flush). \cpp{endl} asegura que el buffer
se vuelca inmediatamente. Para aprendizaje, ambos son equivalentes.
\end{notaprofesor}

% =============================================================================
\section{Material de la sesión}
% =============================================================================

\begin{itemize}
  \item \textbf{Cuaderno}: \texttt{sesiones/sesion02\_introduccion\_cpp.ipynb}
  \item \textbf{Ejemplos}: \texttt{materialInicial/sesion02/main.cpp},
        \texttt{main3.cpp}
  \item \textbf{Notas alumno}: \texttt{notas\_alumno/pdf/sesion02\_introduccion\_cpp.pdf}
\end{itemize}

% =============================================================================
\section{Jueces en línea — Facilitación}
% =============================================================================

\begin{notaprofesor}[Cuándo introducir este tema]
  Después de la demostración de compilación con \texttt{g++}. El juez en línea
  es la extensión natural: ``ya saben compilar localmente; ahora vean cómo
  funciona cuando lo hace un servidor automáticamente''.
\end{notaprofesor}

\subsection{Explicación del concepto (5 min)}

Proyecta omegaUp y muestra un problema sencillo (por ejemplo ``Hello World''
o la suma de dos enteros). Recorre en vivo:

\begin{enumerate}
  \item El enunciado: entrada, salida, restricciones, ejemplos.
  \item El campo de envío: selección de lenguaje, editor o carga de archivo.
  \item El veredicto: muestra qué significa cada código (AC, WA, TLE, CE, RE).
\end{enumerate}

\begin{notaprofesor}
  Haz un envío en vivo con código correcto (AC) y luego uno incorrecto (WA
  o CE). Ver el ciclo completo en tiempo real vale más que explicarlo.
  Si hay acceso a internet limitado, una captura de pantalla del flujo
  funciona igual.
\end{notaprofesor}

\subsection{omegaUp — Por qué usarlo en clase}

\begin{itemize}
  \item \textbf{Mexicano:} fundado en 2011, juez oficial de la OMI. Conectar
        esto con orgullo local suele motivar más de lo esperado.
  \item \textbf{En español:} los enunciados y la interfaz están en español,
        eliminando una barrera de entrada.
  \item \textbf{Grupos y concursos:} puedes crear un grupo de clase y asignar
        problemas directamente. Los alumnos ven su progreso; tú ves quién
        resolvió qué y en cuántos intentos.
  \item \textbf{Estadísticas:} omegaUp registra intentos, tiempo y distribución
        de veredictos, útil para identificar alumnos que necesitan apoyo.
\end{itemize}

\begin{notaprofesor}[Cómo crear un grupo en omegaUp]
  Perfil $\to$ Cursos $\to$ Crear curso. Agrega a los alumnos por su
  nombre de usuario o comparte el enlace de invitación. Desde el curso
  puedes asignar problemas con fecha límite y ver el scoreboard en tiempo
  real durante las sesiones de práctica.
\end{notaprofesor}

\subsection{VJudge — Uso para práctica con problemas internacionales}

VJudge permite crear \textit{contests} con problemas de Codeforces, UVa,
SPOJ, etc. Es útil cuando quieras usar problemas clásicos que no están en
omegaUp. La curva de entrada es un poco mayor; por eso se pide la cuenta
ahora pero se usa activamente más adelante en el curso.

\subsection{Tarea: registro de cuentas}

Pide a los alumnos que registren su nombre de usuario de ambas plataformas
en un formulario (Google Form, hoja de cálculo compartida o el propio
Padlet de la sesión anterior). Necesitarás esos usuarios para agregarlos
al grupo de clase en omegaUp.

% =============================================================================
\section{Rúbrica de práctica}
% =============================================================================

La práctica del cuaderno pide modificar un \cpp{cout} con nombre, carrera y año.

\begin{center}
\begin{tabular}{lcc}
  \toprule
  \textbf{Criterio} & \textbf{Puntos} & \textbf{Evidencia} \\
  \midrule
  Tres \cpp{cout} separados y correctos  & 4 & código ejecutable \\
  Uso correcto de \cpp{endl} o \texttt{\textbackslash n} & 3 & salto de línea visible \\
  Sin errores de compilación             & 3 & celda ejecuta sin error \\
  \bottomrule
  \textbf{Total}                         & 10 & \\
\end{tabular}
\end{center}

\end{document}
